\documentclass[
	fontsize=12pt,          % default font size 12pt
	paper=a4,               % DIN A4 page format
	numbers=noenddot,       % remove dots behind chapter numbers (e.g. 1.5 not 1.5.)
	listof=totoc,           % add list of figures, tables, etc. to ToC
	listof=entryprefix,     % add entry name to figures, tables, etc.
	listof=nochaptergap,    % no chapter gap for figures, tables, etc.
	bibliography=totoc,     % add bibliography to ToC but without a chapter number
	parskip=half            % half line spacing between paragraphs
	openany                 % chapters can start on any page
]{scrbook}
\usepackage{tcolorbox}

\include{../../for_latex/preamble}
\title{\Huge\textbf{Strecke}}
\author{Luis Staudt}
\date{}

% Code-Highlighting für Java
\definecolor{codegray}{rgb}{0.5,0.5,0.5}
\definecolor{codepurple}{rgb}{0.58,0,0.82}
\definecolor{backcolour}{rgb}{0.95,0.95,0.92}
\definecolor{keywordcolor}{rgb}{0.8,0.47,0.13}

\lstdefinestyle{java}{
	language=,
	basicstyle=\ttfamily\small,
	keywordstyle=\bfseries\color{blue!70!black},
	stringstyle=\color{red!60!black},
	commentstyle=\itshape\color{gray},
	numbers=left,
	numberstyle=\tiny\color{gray},
	numbersep=8pt,
	frame=single,
	framerule=0.4pt,
	rulecolor=\color{gray!50},
	breaklines=true,
	showstringspaces=false,
	tabsize=4,
	xleftmargin=1.5em,
	framexleftmargin=1.5em,
	aboveskip=0.8em,
	belowskip=0.6em
}
\lstset{style=java}

\newtcolorbox{hinweis}[1][Hinweis]{
	colback=blue!5,
	colframe=blue!40!black,
	fonttitle=\bfseries,
	title=#1,
	boxrule=0.5pt,
	arc=2pt
}

\newtcolorbox{beispiel}[1][Beispiel]{
	colback=green!5,
	colframe=green!40!black,
	fonttitle=\bfseries,
	title=#1,
	boxrule=0.5pt,
	arc=2pt
}

\begin{document}

% --- Titelblock ---
\begin{center}
	\rule{\textwidth}{0.8pt}\\[0.6em]
	{\LARGE\bfseries Intensivtutorium}\\[0.3em]
	{\large Programmierung}\\[0.6em]
	\begin{tabular}{r l @{\hspace{2cm}} r l}
		\textbf{Semester:}      & Sommersemester 2026                    & \textbf{Tutor:}         & Luis Staudt \\
		\textbf{Studiengänge:}  & MIB, OMB, PEB, MVB         & & \\
	\end{tabular}\\[0.6em]
	\rule{\textwidth}{0.8pt}
\end{center}
\vspace{1em}

% --- Seitenkopf und -fuß (ab Seite 2) ---
\pagestyle{fancy}
\fancyhf{}
\lhead{\small Intensivtutorium -- Grundlagen der Programmierung}
\rhead{\small SS 2026}
\cfoot{\thepage}
\renewcommand{\headrulewidth}{0.4pt}


% =============================================
% Aufgabe 3 -- Klasse Strecke
% =============================================

\section*{Aufgabe 3: Klasse \texttt{Strecke} \hfill (20 Punkte)}

Gegeben ist die aus Vorlesung und Übung bekannte Klasse \texttt{Punkt} (Attribute \texttt{xKoordinate}, \texttt{yKoordinate} als \texttt{double}, Methode \texttt{berechneAbstandZuPunkt(Punkt p)}).

Implementieren Sie eine Klasse \texttt{Strecke}, die eine gerade Verbindung zwischen zwei Punkten in der Ebene beschreibt.
Eine \texttt{Strecke} wird über ihre beiden Endpunkte (jeweils ein \texttt{Punkt}) erzeugt.

\begin{center}
	\begin{tikzpicture}[>=Stealth, scale=1]
		% Achsen
		\draw[->, line width=0.5pt, gray!60] (-0.4,0) -- (5.2,0) node[right, black] {\small $x$};
		\draw[->, line width=0.5pt, gray!60] (0,-0.4) -- (0,3.6) node[above, black] {\small $y$};

		% Punkte
		\coordinate (B) at (1,3);
		\coordinate (A) at (4,1);

		% Strecke
		\draw[line width=1.1pt, blue!60!black] (B) -- (A);

		% Punkt B (links)
		\fill[blue!60!black] (B) circle (2.4pt);
		\node[above] at (B) {\small $(1\,|\,3)$};

		% Punkt A (rechts)
		\fill[blue!60!black] (A) circle (2.4pt);
		\node[above right] at (A) {\small $(4\,|\,1)$};

		% Anfang/Ende Beschriftung
		\node[below left, blue!60!black] at (B) {\footnotesize Anfang};
		\node[below right, blue!60!black] at (A) {\footnotesize Ende};
	\end{tikzpicture}
\end{center}

Die Klasse \texttt{Strecke} soll Folgendes leisten:

\begin{enumerate}[label=\textbf{\arabic*.)}]
	\item Eine \texttt{Strecke} wird aus zwei \texttt{Punkt}-Objekten erzeugt.
		Dabei soll sie ihre Endpunkte \textbf{stets von links nach rechts} verwalten: Der Endpunkt mit der \textbf{kleineren x-Koordinate} ist der Anfang der Strecke, der andere das Ende.
		Bei gleicher x-Koordinate bleibt die übergebene Reihenfolge erhalten.

	\item Sie kann ihre \textbf{Länge} zurückgeben.

	\item Sie kann ihre \textbf{Steigung} zurückgeben:
		\[
			m = \frac{y_{ende} - y_{anfang}}{x_{ende} - x_{anfang}}
		\]
\end{enumerate}

\begin{beispiel}
\begin{lstlisting}
Punkt a = new Punkt();  a.xKoordinate = 4;  a.yKoordinate = 1;
Punkt b = new Punkt();  b.xKoordinate = 1;  b.yKoordinate = 3;

Strecke s = new Strecke(a, b);              // b liegt links -> Anfang = b
System.out.println(s.anfang.xKoordinate);   // 1.0
System.out.println("Laenge:   " + s.berechneLaenge());    // 3.6055...
System.out.println("Steigung: " + s.berechneSteigung());  // -0.6667
\end{lstlisting}
\end{beispiel}

\end{document}
